\section{Dictionnaire}

## 🐍 Cheat Sheet : Les Dictionnaires (`dict`)

> **Règle d'or : Clé / Valeur**
> Un dictionnaire ne fonctionne pas avec des index numériques (comme les listes), mais avec un système d'association. Il associe une **clé** (qui doit être unique et immuable, souvent une `str` ou un `int`) à une **valeur** (qui peut être de n'importe quel type). Les dictionnaires sont modifiables (mutables).

### 1. Opérations Fondamentales

| Syntaxe abstraite | Action |
| --- | --- |
| `len([dict])` | Renvoie le nombre total de paires clé/valeur (`int`). |
| `[dict][[clé]]` | Accède (lit) la valeur associée à cette clé précise. **(Provoque une erreur si la clé n'existe pas).** |
| `[dict][[clé]] = [valeur]` | **Ajout / Modification** : Si la clé existe, met à jour sa valeur. Si elle n'existe pas, crée une nouvelle paire clé/valeur. |
| `[clé] in [dict]` | Opérateur d'appartenance : renvoie `True` si la **clé** (et non la valeur) existe dans le dictionnaire. |

### 2. Accès Sécurisé aux Valeurs

| Méthode abstraite | Action |
| --- | --- |
| `[dict].get([clé])` | Renvoie la valeur associée à la clé. Si la clé n'existe pas, renvoie `None` (sans faire planter le programme). |
| `[dict].get([clé], [défaut])` | Renvoie la valeur associée à la clé. Si la clé n'existe pas, renvoie la valeur de `[défaut]` spécifiée. |

### 3. Vues et Itérations (Très utile pour les boucles `for`)

Ces méthodes renvoient des objets itérables dynamiques ("vues") représentant le contenu du dictionnaire.

| Méthode abstraite | Action |
| --- | --- |
| `[dict].keys()` | Renvoie une vue contenant uniquement toutes les **clés**. |
| `[dict].values()` | Renvoie une vue contenant uniquement toutes les **valeurs**. |
| `[dict].items()` | Renvoie une vue contenant toutes les paires sous forme de tuples : `(clé, valeur)`. |

### 4. Suppression

| Syntaxe abstraite | Action |
| --- | --- |
| `del [dict][[clé]]` | Instruction native qui supprime définitivement la paire associée à cette clé. (Erreur si clé absente). |
| `[dict].pop([clé])` | Supprime la paire associée à cette clé **et renvoie** la valeur supprimée. |
| `[dict].pop([clé], [défaut])` | Identique à `.pop()`, mais renvoie la valeur `[défaut]` si la clé n'existe pas (évite l'erreur). |
| `[dict].popitem()` | Supprime et renvoie la dernière paire insérée dans le dictionnaire sous forme de tuple `(clé, valeur)`. |
| `[dict].clear()` | Supprime la totalité des paires (le dictionnaire devient `{}`). |

### 5. Fusion et Mise à jour massive

| Méthode abstraite | Action |
| --- | --- |
| `[dict].update([autre_dict])` | Modifie le dictionnaire sur place en y ajoutant toutes les paires de `[autre_dict]`. Si des clés sont communes, les valeurs sont écrasées par celles du nouveau dictionnaire. |